\chapter{Metrik Fonksiyonlar ve Sözleşmeler}

Bu bölümde metrik uzaylar arasındaki fonksiyon türleri (izometri, Lipschitz, sürekli)
ve bunların sınıflandırılması incelenmektedir.

%-------------------------------------------------------------------
\section{Konu}

\subsection{Metrik Fonksiyon Sınıfları}

$f: (M_1, d_1) \to (M_2, d_2)$ fonksiyonu:

\begin{table}[h]
\centering
\begin{tabular}{ll}
\toprule
\textbf{Sınıf} & \textbf{Koşul} \\
\midrule
Genişlemez (non-expansive) & $d_2(f(x),f(y)) \leq d_1(x,y)$ \\
Lipschitz & $\exists K>0:\; d_2(f(x),f(y)) \leq K \cdot d_1(x,y)$ \\
Üniform sürekli & $\forall\varepsilon>0\;\exists\delta>0:\; d_1(x,y)<\delta \Rightarrow d_2(f(x),f(y))<\varepsilon$ \\
Sürekli & Her noktada $\varepsilon$-$\delta$ sürekliliği \\
İzometri & $d_2(f(x),f(y)) = d_1(x,y)$ \\
Benzerlik & $\exists c>0:\; d_2(f(x),f(y)) = c \cdot d_1(x,y)$ \\
Homeomorfizma & Bijektif sürekli; ters de sürekli \\
\bottomrule
\end{tabular}
\caption{Metrik fonksiyon sınıfları}
\end{table}

\textbf{Sıralama:} İzometri $\Rightarrow$ Genişlemez $\Rightarrow$ Lipschitz
$\Rightarrow$ Üniform sürekli $\Rightarrow$ Sürekli.

\subsection{Lipschitz Sabiti}

\[
K = \sup_{x \neq y} \frac{d_2(f(x),f(y))}{d_1(x,y)}.
\]

İzometri: $K=1$ ve eşitlik; benzerlik: $K=c$ sabit.

%-------------------------------------------------------------------
\section{Teoremler}

\begin{theorem}
İzometri $\Rightarrow$ homeomorfizma.
\end{theorem}

\begin{theorem}
Benzerlik, $c=1 \Rightarrow$ İzometri.
\end{theorem}

\begin{theorem}
Lipschitz $\Rightarrow$ Üniform sürekli $\Rightarrow$ Sürekli.
\begin{proof}[Kanıt İskeleti]
$\delta = \varepsilon / K$ seçimi Lipschitz $\Rightarrow$ ünif.\ sürekli gösterir.
\end{proof}
\end{theorem}

\begin{theorem}
Kompakt metrik uzaydan metrik uzaya her sürekli fonksiyon üniform süreklidir.
\end{theorem}

%-------------------------------------------------------------------
\section{Algoritmalar}

\subsection{classify\_finite\_metric\_map — $O(|X|^2)$}

\begin{algorithm}
\caption{ClassifyMap$(X, d_1, Y, d_2, f)$}
\begin{algorithmic}[1]
\State $K \leftarrow 0$; isometry $\leftarrow$ True; similarity\_ratio $\leftarrow$ None
\For{each pair $(x_1, x_2)$ with $x_1 \neq x_2$}
    \State ratio $\leftarrow d_2(f(x_1),f(x_2)) / d_1(x_1,x_2)$
    \State $K \leftarrow \max(K, \text{ratio})$
    \If{$d_2(f(x_1),f(x_2)) \neq d_1(x_1,x_2)$}
        \State isometry $\leftarrow$ False
    \EndIf
    \If{similarity\_ratio is None}
        \State similarity\_ratio $\leftarrow$ ratio
    \ElsIf{ratio $\neq$ similarity\_ratio}
        \State similarity\_ratio $\leftarrow$ None
    \EndIf
\EndFor
\State \Return MetricMapProfile(lipschitz\_constant=$K$, isometry=isometry, $\ldots$)
\end{algorithmic}
\end{algorithm}

Karmaşıklık: $O(|X|^2)$.

%-------------------------------------------------------------------
\section{pytop API}

\begin{lstlisting}[language=Python]
from pytop.metric_spaces import FiniteMetricSpace
from pytop.metric_map_taxonomy import (
    classify_finite_metric_map,
    MetricMapProfile,
)
\end{lstlisting}

\texttt{classify\_finite\_metric\_map(fms1, fms2, f)} $\to$ \texttt{MetricMapProfile} döner (Result değil).

\texttt{MetricMapProfile} alanları: \texttt{isometry}, \texttt{similarity}, \texttt{similarity\_ratio},
\texttt{lipschitz\_constant}, \texttt{non\_expansive}, \texttt{bijective}, \texttt{homeomorphism}.

%-------------------------------------------------------------------
\section{Örnekler}

\subsection*{Örnek 5.1 — Özdeşlik: İzometri}

\begin{lstlisting}[language=Python]
f_id = {1: 1, 2: 2, 3: 3, 4: 4}
prof_id = classify_finite_metric_map(fms, fms, f_id)
print("isometry:", prof_id.isometry)
print("lipschitz_constant:", prof_id.lipschitz_constant)
\end{lstlisting}

\begin{verbatim}
isometry: True
lipschitz_constant: 1.0
\end{verbatim}

\subsection*{Örnek 5.2 — 2× Ölçekleme: Benzerlik}

\begin{lstlisting}[language=Python]
f_scale = {1: 2, 2: 4, 3: 6, 4: 8}
prof_scale = classify_finite_metric_map(fms_line1, fms_line2, f_scale)
print("similarity:", prof_scale.similarity)
print("similarity_ratio:", prof_scale.similarity_ratio)
print("lipschitz_constant:", prof_scale.lipschitz_constant)
\end{lstlisting}

\begin{verbatim}
similarity: True
similarity_ratio: 2.0
lipschitz_constant: 2.0
\end{verbatim}

\subsection*{Örnek 5.3 — Sabit Fonksiyon: K=0}

\begin{lstlisting}[language=Python]
f_const = {1: 1, 2: 1, 3: 1, 4: 1}
prof_const = classify_finite_metric_map(fms, fms, f_const)
print("lipschitz_constant:", prof_const.lipschitz_constant)
print("bijective:", prof_const.bijective)
\end{lstlisting}

\begin{verbatim}
lipschitz_constant: 0.0
bijective: False
\end{verbatim}

%-------------------------------------------------------------------
\section{Alıştırmalar}

\subsection*{Kodlama}

\begin{enumerate}[label=K\arabic*.]
\item $\{1,2,3\}$ üzerinde $f(1)=2, f(2)=3, f(3)=1$ (döngü) fonksiyonunu ayrık
      metrikte \texttt{classify\_finite\_metric\_map} ile sınıflandırın.
\item Öklid metrikli $\{0,1,2,3,4\}$ üzerinde $f(x) = x+1$ kaydırma fonksiyonunun
      Lipschitz sabitini hesaplayın.
\item Sabit fonksiyon $f(x) = 1$ için $K=0$ olduğunu doğrulayın.
\end{enumerate}

\subsection*{Teori}

\begin{enumerate}[label=T\arabic*.]
\item İzometri $\Rightarrow$ genişlemez $\Rightarrow$ Lipschitz ($K=1$) zincirini ispatlayın.
\item Lipschitz $\Rightarrow$ üniform sürekli olduğunu $\delta = \varepsilon/K$ seçimi
      ile ispatlayın.
\end{enumerate}
